\section{Introduction}
\label{sec:introduction}

Tax-benefit microsimulation models let researchers and policymakers evaluate
the distributional and fiscal effects of policy changes before they are
enacted, by applying rules to representative individual- or household-level
records rather than to aggregates \citep{bourguignon2006microsimulation}. In
the United Kingdom, the incumbent academic model is \ukmod{}, itself the UK
country instance of \euromod{}, the tax-benefit model of the European
Union \citep{sutherland2013euromod}. \ukmod{} and \euromod{} are documented
in a well-established citation lineage that this paper's structure follows;
so is \textsc{TAXSIM}, the equivalent US individual income tax calculator
\citep{feenberg1993taxsim}. \policyengine{} UK
(\texttt{\href{https://github.com/PolicyEngine/policyengine-uk}{policyengine-uk}})
is an open-source model of the full UK personal tax and benefit system that
has operated as a public, continuously deployed tool since \todo{launch year
-- confirm from repository history or an early release/announcement} but has
not until now had a citable, versioned description of its scope, data
construction, and validation. This paper is that description.

\todo{One paragraph situating \policyengine{} UK relative to \ukmod{} and
official HMRC/DWP costing tools (e.g. IPPR tax-benefit model, if a citable
comparison exists) -- strictly factual, no claims beyond what is
independently verifiable, per the neutral-framing requirement on this paper
(PolicyEngine/policyengine-uk-paper\#1). This paragraph should describe
what \policyengine{} UK models and cites (scope, coverage, data sources);
it should not frame \ukmod{}'s longer citation history as reflecting
better or worse on either model -- the preceding paragraph's citable-paper
timeline is a fact about \policyengine{} UK's own publication history, not
a comparative claim, and this paragraph should not convert it into one.}

\subsection{Versioning}

Software-backed microsimulation models change continuously: parameters are
uprated annually, new programs are added, and the underlying simulation
engine is occasionally replaced. A paper that describes ``the model'' at a
single point in time becomes silently inaccurate as the software evolves.
Following the versioning framing used for the sibling paper on
\populace{}'s imputation operator \citep{populace2026}, this paper is
explicitly versioned: \textbf{v1} documents \policyengine{} UK as deployed in
2026, on the \texttt{policyengine\_core} engine and the \populace{} data
stack described in Section~\ref{sec:data}. A future migration of the
underlying engine or a substantial change in data construction will produce a
v2 revision of this paper rather than retroactively invalidating this one.

\subsection{Scope and contributions}

This paper:

\begin{enumerate}
    \item Documents the tax and benefit programs \policyengine{} UK models,
    their administrative sourcing, and known coverage gaps
    (Section~\ref{sec:model}).
    \item Documents the survey microdata the model runs on -- the \frs{},
    its enhancement with \hmrc{} \spi{} income data, and calibration to
    administrative totals -- and the versioned data pipeline that produces
    it (Section~\ref{sec:data}).
    \item Validates modelled aggregates against an independent
    microsimulation model (\ukmod{}) and official administrative and
    forecast aggregates (\hmrc{}, \dwp{}, \obr{}), and validates
    individual-level calculations against \hmrc{}'s \spi{} microdata
    (Section~\ref{sec:validation}).
    \item States what is not yet validated and where the model's coverage is
    partial, so that users can assess fitness for a given application
    (Section~\ref{sec:discussion}).
\end{enumerate}

Consistent with the scope of this v1, we do not report new simulation
results: every quantity in Section~\ref{sec:validation} is drawn from an
existing, committed validation artifact in the \texttt{policyengine-uk} or
\texttt{populace} repositories, cited to its source file. Where a claim in
this outline does not yet have a cited source, it is marked
\todo{example placeholder}.
